\SourceLeaf{55}$p$ incongruos) habebit. Hinc perspicitur, expressionem ${}^{\delta}\!\surd {1}$ etiam $\delta$ valores diversos habere, quorum indices cum ante allatis prorsus conveniant. Quocirca expressio ${}^{\delta}\!\surd {1}\,(\operatorname{mod.}p)$ huic ${}^{n}\!\surd {1}\,(\operatorname{mod.}p)$ omnino aequivalet, i. e. congruentia $x^\delta\equiv1\,(\operatorname{mod.}p)$ easdem radices habet quas haec, $x^n\equiv1\,(\operatorname{mod.}p)$. Prior autem inferioris erit gradus siquidem $\delta$ et $n$ sunt inaequales.

\emph{Ex.} ${}^{15}\!\surd {1}\,(\operatorname{mod.}19)$ tres habet valores, quia $3$ maxima numerorum $15,18$ mensura communis, hique simul erunt valores expressionis ${}^{3}\!\surd {1}\,(\operatorname{mod.}19)$. Sunt autem hi $1,7,11$.

\Article{62}

Per hanc igitur reductionem id lucramur ut alias congruentias formae $x^n\equiv1$ solvere non sit opus, quam ubi $n$ numeri $p-1$ est divisor. Infra vero ostendemus, congruentias huius formae semper ulterius adhuc deprimi posse, licet praecedentia ad hoc non sufficiant. Unum tamen casum iam hic absolvere possumus scilicet ubi $n=2$. Manifesto enim valores expressionis ${}^{2}\!\surd {1}$ erunt $+1$ et $-1$ quia plures quam duos habere nequit, hique $+1$ et $-1$ semper sunt incongrui nisi modulus sit $=2$, in quo casu ${}^{2}\!\surd {1}$ unum tantum valorem habere posse, per se clarum. Hinc sequitur, $+1$ et $-1$ etiam fore valores expressionis ${}^{2m}\!\surd {1}$ quando $m$ ad $\frac{p-1}{2}$ sit primus. Hoc semper eveniet, quoties modulus est eius indolis ut $\frac{p-1}{2}$ fiat numerus absolute primus (nisi forte $p-1=2m$ in quo casu omnes numeri $1,2,3\GaussFourDots p-1$ sunt radices) \emph{ex. gr.} quando $p=3,5,7,11,23,47,59,83,107$ etc. Tamquam corollarium hic annotetur, indicem ipsius $-1$ semper esse $\equiv\frac{p-1}{2}\,(\operatorname{mod.}p-1)$, quaecunque radix primitiva pro basi accipiatur. Namque $2\operatorname{Ind.}(-1)\equiv0\,(\operatorname{mod.}p-1)$. Quare $\operatorname{Ind.}(-1)$ erit vel $\equiv0$, vel $\equiv\frac{p-1}{2}\,(\operatorname{mod.}p-1)$: $0$ vero semper index ipsius $+1$, atque $+1$ et $-1$ semper indices diversos habere debent (praeter casum $p=2$ ad quem hic respicere operae non est pretium).

\Article{63}

Ostendimus art. 60 expressionem ${}^{n}\!\surd {A}\,(\operatorname{mod.}p)$ habere $\delta$ valores diversos, aut omnino nullum, si fuerit $\delta$ divisor communis maximus numerorum $n,p-1$. Iam uti modo docuimus ${}^{n}\!\surd {A}$ et ${}^{\delta}\!\surd {A}$ aequivalentes esse, si fuerit $A\equiv1$, generalius probabimus, expressionem ${}^{n}\!\surd {A}$ semper ad aliam ${}^{\delta}\!\surd {B}$ reduci posse cui aequivaleat. Illius enim valore quocunque denotato per $x$ erit $x^n\equiv A$; iam
\SourceLeaf{56}sit $t$ valor quicunque expressionis $\frac{\delta}{n}\,(\operatorname{mod.}p-1)$, quam valores reales habere ex art. 31 perspicuum; eritque $x^{tn}\equiv A^t$ at $x^{tn}\equiv x^\delta$ propter $tn\equiv\delta\,(\operatorname{mod.}p-1)$. Quare $x^\delta\equiv A^t$ adeoque quicunque ipsius ${}^{n}\!\surd {A}$ valor erit etiam valor ipsius ${}^{\delta}\!\surd {A^t}$. Quoties igitur ${}^{n}\!\surd {A}$ valores reales habet, expressioni ${}^{\delta}\!\surd {A^t}$ prorsus aequivalens erit, quoniam illa neque alios habet quam haec neque pauciores, licet quando ${}^{n}\!\surd {A}$ nullum valorem realem habet, fieri tamen possit ut ${}^{\delta}\!\surd {A^t}$ valores reales habeat.

\emph{Ex.} Si valores expressionis ${}^{21}\!\surd {2}\,(\operatorname{mod.}31)$ quaeruntur, erit numerorum $21$ et $30$ divisor communis maximus $3$, expressionisque $\frac{3}{21}\,(\operatorname{mod.}30)$ valor aliquis $3$, quare si ${}^{21}\!\surd {2}$ valores reales habet, huic expressioni ${}^{3}\!\surd {2^3}$ sive ${}^{3}\!\surd {8}$ aequivalebit, invenieturque revera, posterioris expressionis valores qui sunt $2,10.19$ etiam priori satisfacere.

\Article{64}

Ne autem hanc operationem incassum suscepisse periclitemur, regulam investigare oportet, per quam statim diiudicari possit utrum ${}^{n}\!\surd {A}$ valores reales admittat necne. Quodsi tabula indicum habetur, res in promtu est; namque ex art. 60 manifestum est, valores reales dari, si ipsius $A$ index, radice quacunque primitiva pro basi accepta. per $\delta$ sit divisibilis, sin vero minus, non dari. Attamen hoc etiam absque tali tabula inveniri potest. Posito enim indice ipsius $A=k$, si hic fuerit per $\delta$ divisibilis, erit $\frac{k(p-1)}{\delta}$ per $p-1$ divisibilis et vice versa. Atqui numeri $A^{\frac{p-1}{\delta}}$ index erit $\frac{k(p-1)}{\delta}$. Quare si ${}^{n}\!\surd {A}\,(\operatorname{mod.}p)$ habet valores reales, $A^{\frac{p-1}{\delta}}$ unitati congruus erit, sin minus, incongruus. Ita in exemplo art. praec. habetur $2^{10}=1024\equiv1\,(\operatorname{mod.}31)$, unde concluditur ${}^{21}\!\surd {2}\,(\operatorname{mod.}31)$ valores reales habere.
Similiter certiores hinc fimus, ${}^{2}\!\surd {-1}\,(\operatorname{mod.}p)$ semper valores binos reales habere, quando $p$ sit formae $4m+1$, nullum vero, quando $p$ sit formae $4m+3$; propter $(-1)^{2m}=1$ et $(-1)^{2m+1}=-1$. Elegans hoc theorema, quod vulgo ita profertur: \emph{Si $p$ est numerus primus formae $4m+1$, inveniri potest quadratum $aa$, ita ut $aa+1$ per $p$ fiat divisibilis; si vero $p$ est formae $4m-1$, tale quadratum non datur}, hoc modo demonstratum est ab ill. Eulero, \emph{Comm. nov. Acad. Petrop. T.} XVIII \emph{p.} 112 ad annum 1773. Demonstrationem aliam iam multo ante dederat, \emph{Comm. nov. T.} V \emph{p.} 5 qui prodiit a. 1760. In dissert. priori, \emph{Comm. nov. T.} IV \emph{p.} 25, rem nondum
\SourceLeaf{57}perfecerat. Postea etiam ill. La Grange theorematis demonstrationem tradidit, \emph{Nouveaux Mém. de l'Ac. de Berlin A.} 1775 \emph{p.} 342. Aliam adhuc demonstrationem in sectione sequenti ubi proprie de hoc argumento agendum erit, dabimus.

\Article{65}

Postquam omnes expressiones ${}^{n}\!\surd {A}\,(\operatorname{mod.}p)$ ad tales reducere docuimus, ubi $n$ divisor numeri $p-1$, criteriumque nacti sumus utrum valores reales admittat, necne, tales expressiones ${}^{n}\!\surd {A}\,(\operatorname{mod.}p)$ ubi $n$ ipsius $p-1$ est divisor accuratius considerabimus. Primo ostendemus, quam relationem valores singuli expressionis inter se habeant, tum artificia quaedam trademus, quorum auxilio unus valor expressionis saepenumero inveniri possit.

\emph{Primo,} quando $A\equiv1$, atque $r$ aliquis ex $n$ valoribus expressionis ${}^{n}\!\surd {1}\,(\operatorname{mod.}p)$, sive $r^n\equiv1\,(\operatorname{mod.}p)$, omnes etiam ipsius $r$ potestates erunt valores istius expressionis; horum autem totidem erunt diversi quot unitates habet exponens ad quem $r$ pertinet (art. 48). Quodsi igitur $r$ est valor ad exponentem $n$ pertinens, potestates ipsius $r$ hae $r,r^2,r^3\GaussFourDots r^n$ (ubi loco ultimae \emph{unitas} substitui potest) omnes expressionis ${}^{n}\!\surd {1}\,(\operatorname{mod.}p)$ valores involvent. Qualia autem subsidia exstent ad tales valores inveniendos qui ad exponentem $n$ pertineant, in Sect. VIII fusius explicabimus.

\emph{Secundo.} Quando $A$ unitati est incongruus, unusque valor expressionis ${}^{n}\!\surd {A}\,(\operatorname{mod.}p)$ notus, qui sit $z$, reliqui hoc modo inde deducuntur. Sint valores expressionis ${}^{n}\!\surd {1}$ hi

\[1,\ r,\ r^2\ldots r^{n-1}\]

(uti modo ostendimus), eruntque omnes expr. ${}^{n}\!\surd {A}$ valores hi

\[z,\ zr,\ zr^2\GaussFourDots zr^{n-1}\]

namque omnes hos congruentiae $x^n\equiv A$ satisfacere inde manifestum quod, posito quocunque eorum $\equiv zr^k$, potestas ipsius $n^{\mathrm{ta}}$, $z^n r^{nk}$, propter $r^n\equiv1$ et $z^n\equiv A$, ipsi $A$ fit congrua: omnes diversos esse ex art. 23 facile intelligitur; plures autem valores quam hos quorum numerus est $n$, expressio ${}^{n}\!\surd {A}$ habere nequit. Ita \emph{ex. gr.} si alter expressionis ${}^{2}\!\surd {A}$ valor est $z$, alter erit $-z$. Deni\SourceLeaf{58}que hinc concludendum omnes valores expr. ${}^{n}\!\surd {A}$ inveniri non posse, nisi simul omnes valores expr. ${}^{n}\!\surd {1}$ constent.

\Article{66}

Secundum quod nobis proposueramus fuit docere, in quo casu unus expressionis ${}^{n}\!\surd {A}\,(\operatorname{mod.}p)$ valor (ubi $n$ supponitur esse divisor ipsius $p-1$) directe inveniri possit. Hoc evenit quando aliquis valor potestati alicui ipsius $A$ congruus evadit, qui casus quum haud raro occurrat, aliquantum huic rei immorari non superfluum erit. Sit talis valor, \emph{si quis datur} $z$, sive $z\equiv A^k$ et $A\equiv z^n\,(\operatorname{mod.}p)$. Hinc colligitur $A\equiv A^{kn}$; quare si numerus $k$ habetur, ita ut sit $A\equiv A^{kn}$, $A^k$ erit valor quaesitus. At huic conditioni aequivalet ista, ut sit $1\equiv kn\,(\operatorname{mod.}t)$, designante $t$ exponentem ad quem pertinet $A$ (art. 46, 48). Ut vero haec congruentia possibilis sit, requiritur, ut sit $n$ ad $t$ primus. Hoc in casu erit $k\equiv\frac{1}{n}\,(\operatorname{mod.}t)$; si vero $t$ et $n$ divisorem communem habent, nullus valor $z$ potestati ipsius $A$ congruus esse potest.

\Article{67}

Quum autem ad hanc solutionem ipsum $t$ novisse oporteat, videamus quomodo procedere possimus, si hunc numerum ignoremus. Primo facile intelligitur, $t$ ipsum $\frac{p-1}{n}$ metiri debere, siquidem ${}^{n}\!\surd {A}\,(\operatorname{mod.}p)$ valores reales habeat, uti hic semper supponimus. Sit enim quicunque valor $y$, eritque tum $y^{p-1}\equiv1$, tum $y^n\equiv A\,(\operatorname{mod.}p)$; quare elevando partes posterioris congruentiae ad potestatem $\frac{p-1}{n}^{\mathrm{tam}}$, fiet $A^{\frac{p-1}{n}}\equiv1$; adeoque $\frac{p-1}{n}$ per $t$ divisibilis (art. 48). Iam si $\frac{p-1}{n}$ ad $n$ est primus, congruentia art. praec. $kn\equiv1$ etiam secundum modulum $\frac{p-1}{n}$ solvi poterit, manifestoque valor ipsius $k$ congruentiae secundum modulum hunc satisfaciens eidem etiam secundum modulum $t$, qui ipsum $\frac{p-1}{n}$ metitur, satisfaciet (art. 5). Tum igitur quod quaerebatur inventum. Si vero $\frac{p-1}{n}$ ad $n$ non est primus, omnes ipsius $\frac{p-1}{n}$ factores primi qui simul ipsum $n$ metiuntur ex $\frac{p-1}{n}$ eiiciantur. Hinc nanciscemur numerum $\frac{p-1}{nq}$, ad $n$ primum, designante $q$ productum ex omnibus illis factoribus primis, quos eiecimus. Quodsi iam conditio ad quam in artic. praec. pervenimus ut $t$ ad $n$ sit primus locum habet, $t$ etiam ad $q$ erit primus adeoque etiam ipsum $\frac{p-1}{nq}$ metietur. Quare si congruentia $kn\equiv1\,(\operatorname{mod.}\frac{p-1}{nq})$ solvitur (quod fieri potest quia $n$ ad $\frac{p-1}{nq}$ primus), valor ipsius $k$ etiam secundum modulum
